\section{Task graph}
\label{sec:taskgraph}

This section documents how the HII search is wired into \swift's task
graph, focused on the stellar/radiation part only (not hydro or
gravity, except where radiation depends on them). All line references
are current as of \texttt{darwin/gear\_radiation},
2026-07-16.

\subsection{Why this is not a regular self/pair task}
\label{sec:not-regular}

An ordinary \swift self/pair loop (e.g.\ hydro density) is created once
per pair of neighbouring cells at \texttt{hydro.super}, and
\texttt{scheduler\_splittasks} recursively breaks each one into many
independent, finer self/pair tasks wherever
\texttt{cell\_can\_split\_self/pair\_hydro\_task} says the cut-off
radius still fits inside a sub-cell -- so the actual density/force
calculation is spread across potentially hundreds of small,
independently-schedulable tasks that each touch only two specific
cells.

The radiation search does not work this way, for two reasons: (1) the
relevant length scale, $h_{\rm hii}$, is not a fixed smoothing length
but a search radius that grows dynamically over many rebuild cycles,
so no single tree level is guaranteed to contain it; and (2) the
consumption logic is inherently sequential and global to one star (a
star's photon budget must be spent in strict distance order across
\emph{all} its candidates, not independently per cell-pair). So the
implementation splits the work into two different kinds of task:

\begin{itemize}
  \item \textbf{\texttt{task\_subtype\_stars\_radiation\_in} /
    \texttt{\_out}} (self and pair): created once at the
    \emph{top-level} cell grid (27-neighbour stencil,
    \texttt{engine\_make\_radiationloop\_tasks\_mapper},
    \texttt{engine\_maketasks.c:3902}), \textbf{as implicit tasks}
    (\texttt{implicit=1} in \texttt{scheduler\_addtask}) -- they run no
    compute code at all (no case for them in \texttt{runner\_main.c}'s
    dispatch switch). Their only job is to exist as dependency nodes:
    every drift/sort/cooling task that must finish before the search
    reads that region unlocks the matching \texttt{\_in} task; the
    \texttt{\_out} task unlocks \texttt{stars\_out} once the search is
    done. They are never split further by
    \texttt{scheduler\_splittasks}.
  \item \textbf{\texttt{task\_type\_stars\_hii\_ionization\_feedback}}
    (subtype \texttt{none}): a \emph{single, coarse-grained} task
    created once per \texttt{radiation\_level} cell
    (\texttt{engine\_make\_hierarchical\_tasks\_radiation\_subgrid},
    \texttt{engine\_maketasks.c:2109}). This is the only task that
    actually runs code (\texttt{runner\_do\_stars\_hii\_\allowbreak
    ionization\_feedback}, dispatched at \texttt{runner\_main.c:582}).
    It performs its own manual, run-time recursion down to the
    ``Working Level'' and its own per-star neighbour search internally
    -- see Section~\ref{sec:runtime}. From the scheduler's point of
    view this is one big task per region, not many small ones.
\end{itemize}

\subsection{\texttt{radiation\_level}: a coarsening criterion decoupled from \texttt{hydro.super}}

\texttt{cell\_can\_split\_pair\_radiation\_subgrid\_task} /
\texttt{cell\_can\_split\_self\_radiation\_subgrid\_task}
(\texttt{cell.h:1201}, \texttt{1239}) are radiation's own splitting
criteria, structurally identical to the hydro ones but with an extra
term: a cell stops splitting (and \texttt{radiation\_level} is pinned
there) once
\begin{equation}
  \text{space\_stretch} \times \gamma_{\rm kernel} \times h_{{\rm hii},\max} \ge 0.5\, d_{\min},
\end{equation}
\emph{in addition to} the ordinary hydro/star/sink/BH $h_{\max}$
terms, and also \texttt{c->hydro.super == NULL} (i.e.\ radiation never
splits finer than \texttt{hydro.super}, only coarser). Since $h_{\rm
hii}$ can grow far beyond an ordinary smoothing length as a star's
search radius expands, \texttt{radiation\_level} can sit strictly
\emph{above} \texttt{hydro.super} -- one radiation task then covers
several \texttt{hydro.super} subtrees at once. When $h_{\rm hii}$ is
small this term is inert and \texttt{radiation\_level == hydro.super}
for free. \texttt{cell\_set\_super\_radiation\_subgrid}
(\texttt{cell.c:1361}) walks the tree and stamps
\texttt{c->stars.radiation\_level} to the shallowest ancestor holding
a \texttt{radiation\_in} link.

\textbf{Known, code-verified, not-yet-observed risk} (guarded by a
\texttt{SWIFT\_DEBUG\_CHECKS} assertion, \texttt{cell.c:1387}): the
run-time search
(\texttt{runner\_dosub\_stars\_hii\_ionization\_feedback}) only ever
reads \texttt{radiation\_level}'s \emph{own} link list, never descends
into its children's links. If two neighbouring regions end up with
\texttt{hydro.super} at different depths (e.g.\ from non-uniform
density), a cell can hold its own \texttt{radiation\_in} link while an
ancestor is found first and claims \texttt{radiation\_level} -- the
child's link then becomes permanently invisible to the flat walk, a
silent search-incompleteness bug, not a crash. Confirmed by
code-reading only; uniform-density ICs never trigger the depth
mismatch, but the clump examples in this branch are exactly the kind
of non-uniform setup that could. See item~T1 in
Section~\ref{sec:open}.

\subsection{Dependency wiring}

For every \texttt{radiation\_in} self/pair task \texttt{t} covering
cell(s) \texttt{ci}/\texttt{cj}, \texttt{engine\_\allowbreak
make\_extra\_radiationloop\_tasks\_mapper}
(\texttt{engine\_maketasks.c:3608}) creates the matching
\texttt{radiation\_out} task and wires, via
\texttt{engine\_radiation\_wire\_super\_deps}
(\texttt{engine\_maketasks.c:3574}, which recurses down to every
\texttt{hydro.super} beneath a coarse \texttt{radiation\_level} cell,
visiting each disjoint subtree exactly once):
\[
  \underbrace{\{\text{drift}_{\rm hydro},\ \text{sorts}_{\rm hydro},\ \text{cooling\_out},\ \text{drift}_{\rm stars},\ \text{feedback\_ghost}\}}_{\text{per hydro.super beneath } ci/cj}
  \;\to\; t_{\rm in} \;\to\; \text{hii\_ionization\_feedback} \;\to\; t_{\rm out} \;\to\; \text{stars\_out}_{\rm hydro.super},
\]
plus \texttt{t\_out} $\to$ \texttt{timestep\_sync} (once, at the main
\texttt{super}, if the policy is active). A
\texttt{SWIFT\_DEBUG\_CHECKS} invariant enforces that a
\texttt{radiation\_in} pair task always connects two \emph{distinct}
\texttt{radiation\_level} regions (\texttt{engine\_maketasks.c:3691}) --
otherwise both sides would wire identical unlocks onto the same
\texttt{hii\_ionization\_feedback} task, a duplicate-unlock error.

\subsection{Activation (unskip)}

\texttt{cell\_unskip\_radiation\_tasks} (\texttt{cell\_unskip.c:2243})
walks \texttt{c->stars.radiation\_in} and \texttt{radiation\_out}
each timestep. For every active task it also activates
drift/sync/\texttt{stars\_in} (or \texttt{stars\_out}) over
\emph{every} \texttt{hydro.super} beneath the possibly-coarse
\texttt{radiation\_level} cell(s)
(\texttt{cell\_radiation\_activate\_supers\_in/out}, factored out
specifically because a single \texttt{c->hydro.super} dereference
would be NULL or wrong for a coarsened region). The pair branch also
runs the radiation-specific rebuild check,
\texttt{cell\_need\_rebuild\_for\_radiation\_pair} -- the ordinary
\texttt{cell\_need\_rebuild\_for\_stars\_pair} is blind to $h_{\rm
hii}$ growth between rebuilds and would under-trigger.
\texttt{c->stars.hii\_ionization\_feedback} itself is unskipped
unconditionally alongside the other per-cell tasks
(\texttt{cell\_unskip.c:2731}) whenever the cell needs activating.

\subsection{Locking}

\texttt{hii\_ionization\_feedback}'s search reads and writes gas across
\texttt{ci}'s own subtree plus every neighbouring region reachable
through \texttt{ci->stars.radiation\_in} (\texttt{task.c:805}, lock;
\texttt{task.c:\allowbreak 577}, unlock). Locking proceeds in
\texttt{radiation\_in}'s link order: \texttt{ci}'s own star and hydro
trees first, then each linked pair task's ``other'' cell's hydro tree
(locking a region's tree also blocks its descendants/ancestors via
\texttt{cell\_locktree}'s hold-counter, so one lock per neighbouring
region suffices -- no need to enumerate every \texttt{hydro.super}
beneath it individually). If any lock in the chain fails, everything
already acquired for this task is rolled back and the task is
requeued -- same contention pattern as any other two-cell pair lock,
not a novel risk, but link order is only locally (not globally)
consistent across two regions locking each other, so genuine
lock-order contention between neighbours is possible under load.

\subsection{Run-time recursion and neighbour search}
\label{sec:runtime}

Once \texttt{hii\_ionization\_feedback} actually runs
(\texttt{runner\_do\_stars\_hii\_ionization\_feedback},
\texttt{runner\_radiation\_feedback.c:52}), it manually recurses down
\texttt{c->progeny} (mirroring what \texttt{scheduler\_splittasks}
would have done for an ordinary task, but at run time inside a single
task) until it reaches the ``Working Level'' -- a cell no smaller than
needed to contain the current interaction limit -- then hands off to
\texttt{runner\_dosub\_stars\_hii\_ionization\_feedback}
(\texttt{runner\_radiation\_feedback.c:158}). For each active star at
that level:
\begin{enumerate}
  \item Gather candidates within the current search radius via
    \texttt{runner\_do(self/dopair)\_stars\_hii\_ionization\_feedback},
    walking \texttt{c}'s own cell and every cell reachable through
    \texttt{c->stars.radiation\_level->stars.radiation\_in} (the same
    link list used for locking/dependencies) into a fixed-size,
    distance-sorted buffer (\texttt{max\_ngbs}), filtered by
    \texttt{feedback\_part\_can\_be\_ionized} (not already tagged, cold
    enough, dense enough) and, if angular splitting is on, by whether
    the candidate's assigned pixel still has budget.
  \item Consume the star's budget on the sorted buffer, closest first
    (Eq.~\ref{eq:recomb-cost}), applying the deterministic/probabilistic
    boundary-particle rule at exhaustion.
  \item Retry: if the buffer filled up (more candidates may exist at
    the \emph{same} radius, just bumped out by the buffer's size cap),
    retry at the same radius, up to
    \texttt{Stars:max\_HII\_retry\_full\_buffer} times. If the buffer
    did \emph{not} fill up but budget remains, the radius itself is
    the bottleneck -- expand it by
    \texttt{Stars:HII\_radius\_expansion\_factor}, up to
    \texttt{Stars:max\_HII\_radius\_expansion\_tries} times, bounded by
    \texttt{Stars:max\_HII\_search\_radius}. These two retry kinds are
    genuinely different failure modes and must not be conflated (see
    the comment at \texttt{runner\_radiation\_feedback.c:213}); the
    loop breaks once the budget is exhausted or both retry budgets run
    out.
\end{enumerate}
This retry/expansion loop, not \texttt{scheduler\_splittasks}, is what
makes the search radius-adaptive: it is why hitting
\texttt{Stars:max\_HII\_search\_radius} (Section~\ref{sec:validation},
Table~\ref{tab:matrix} rows 2/6) shows up as the search radius simply
never being allowed to expand past the configured ceiling, not as a
scheduler-level task-graph error.
